\documentclass{article}
\usepackage[utf8]{inputenc}
\newcommand{\comment}[1]{}
\usepackage{tabularx}
\usepackage{verbatim}
\usepackage{hyperref}
\usepackage{mathtools}
\usepackage{float}
\usepackage{subcaption}
\usepackage{textcomp}
\usepackage[a4paper, margin={1in}]{geometry}
\usepackage[english]{babel}
\usepackage{graphicx}
\usepackage{amsmath}
\usepackage{amsfonts}
\usepackage{fontspec}
\setmonofont{FreeMono}
\usepackage{minted}
\setminted{breaklines=true}
\newmintinline[lean]{lean4}{bgcolor=white}
\newminted[leancode]{lean4}{fontsize=\footnotesize}
\usemintedstyle{tango}

\newcolumntype{Y}{>{\centering\arraybackslash}X}

\usepackage{amsthm}
\usepackage{authblk}
\usepackage{biblatex}
\usepackage{mathpartir}
\usepackage{stmaryrd}
\usepackage{amssymb}
\newtheorem{lemma}{Lemma}
\newtheorem{theorem}{Theorem}
\addbibresource{./refs.bib}

\newcommand{\set}{\mathrm{Set}}
\newcommand{\sort}[1]{\mathrm{Sort}\ #1}
\newcommand{\affarr}{\rightarrowtail}
\newcommand{\affetaarr}{\rightarrowtail_{\eta}}
\newcommand{\etaarr}{\rightarrow_{\eta}}

\title{Decidable Type-Checking for Extensional Type Theory in Combinator Calculus}
\author[1]{Lexi Aiello}
\author[2]{Eric Klavins}
\affil[1]{Department of Philosophy, University of Washington}
\affil[2]{Electrical \& Computer Engineering Department, University of Washington}

\section{Introduction}

\begin{comment}

  High level flow:
  1. Abstract
  2. Introduction
    - Mention features
    - Describe results, and how they achieve the features
    2a. Background yadda yadda, segue into
  3. A recipe for computable function extensionality (overall description)
    3a. Affine insight: strong normalization by default (see below)
    3b. Cat. combinator insight
  ?. The categorical combinator insight
  


  Features list:
  - Dependent types
  - Decidable type-checking
  - Computable function extensionality as a metatheorem
  - Computable propositional equality reflection as a metatheorem
  - Fully combinatory with no variables
  - Strong normalization
  - Based on affine logic and categorical combinators
  - Compiler from a fragment of Lean to the calculus, supporting complex dependent Lean functions (e.g., the dependent K combinator written in Lean)

  Results list:
  - Decidable type-checking for the extensional, affine combinator calculus based on categorical combinators
  - Strong normalization inherited by default
  - Funext proven as a computable def metatheorem
    - Due to lack of open terms in combinators and connection between cat. combinators and CCC's
    - Equalities between functions can be proven with refl x
    - Note: this is important, because these combination of things seems to produce decidable funext:
      - Affine logic (SN and decidable beta-eq by default)
        - Since SN by default, we can eagerly reduce all types
          - No general conversion constructor -> only one canonical judgment exists -> extensionality is possible
      - Combinatory basis (no open terms)
  - The extensional elimination equality reflection rule also proven as a metatheorem
  - Compiler from the calculus of constructions (i.e., Lean code) to DCC

  - [ ] Results list
  - Extension of categorical combinators to dependent types
  - Affine insight: cat. combinators are strongly normalizing with modification to the pair rule
    - Consequences of this:
      - Size measure gives strong normalization
        - Due to SN, we have UNF.
        - We separately show confluence.
        - We have global confluence due to wellfoundedness by size.
        - Decidable beta-eq
          - IMPORTANT: No need for a general fallback conversion constructor in the judgment relation
            - This makes metatheory easier. Avoid hairy induction
            - This means there is only one canonical judgment for a given term
  - Cat. combinator insight: eta-reduction in the categorical combinators corresponds to the primitive morphisms of CCC's
    - Unlike in lambda-calculus, there are no open terms
      - Eta-canonical forms are well-behaved
      - We use a syntactic eta-reduction approach to observe the behavior of functions on ``holes'', whose shapes are determined by the function's domain
      - This is particularly fitting and feels quite natural to cat. combinators
        - There only exist functions of arity 1 or 2. Curry is an example of a function of arity 1.
          - Two holes is sufficient to test behavior, since Curry is the only non-unary function
      - IMPORTANT: We have proven funext as a metatheorem for eta-reduction
        - -> show some examples
        - This is possible since:
          - We already have SN, so we can reduce fully
  
  - reference observational type theory for comparison, since we also have decidable funext
  - example of using extensionality in proof
  - 

Things I want to cover

\end{comment}
